\documentclass[a4paper,12pt]{article}
\usepackage[english]{babel}
\usepackage[utf8]{inputenc}
\usepackage[T1]{fontenc}
\usepackage{enumitem}
\usepackage{amsmath,amssymb}
\usepackage{geometry}
\geometry{left=25mm,right=25mm,top=25mm,bottom=25mm}
\usepackage{parskip}

\newif\ifshowsolutions
\showsolutionstrue  % comment to hide solutions

\begin{document}

\begin{center}
  {\LARGE\bfseries Control Technology \\[6pt] Exercises}\\[12pt]
  \rule{\textwidth}{0.4pt}
\end{center}


\section{Tank Level Control System}

A liquid tank has cross-sectional area $A$. The inflow is $q_{in}(t)$, outflow is proportional to level $h(t)$.

\[
A\frac{dh(t)}{dt} = q_{in}(t) - \frac{1}{R}h(t)
\]

Parameters:

\[
A=5,\quad R=2
\]

\begin{enumerate}

\item[(a)] Derive the transfer function $G(s)=H(s)/Q_{in}(s)$.

\item[(b)] Determine system type and stability.

\item[(c)] Calculate the step response:
\begin{enumerate}
\item[(i)] steady-state gain
\item[(ii)] time constant
\end{enumerate}

\item[(d)] Interpret the frequency response.

\item[(e)] PI control design:
\[
C(s)=K_P+\frac{K_I}{s}
\]
\begin{enumerate}
\item[(i)] Derive the closed-loop characteristic equation.
\item[(ii)] Determine the steady-state error for a step input.
\end{enumerate}

\item[(f)] Engineering interpretation:
Why is PI control almost always used for tank systems?

\end{enumerate}

\ifshowsolutions
Solution:

\begin{enumerate}

\item[(a)]~\\

Laplace:

\[
AsH(s)=Q_{in}(s)-\frac{1}{R}H(s)
\]

\[
H(s)\left(As+\frac{1}{R}\right)=Q_{in}(s)
\]

\[
G(s)=\frac{H(s)}{Q_{in}(s)}=\frac{1}{As+\frac{1}{R}}
\]

Insert values:

\[
G(s)=\frac{1}{5s+0.5}
\]

\item[(b)]~\\

First order system → stable (pole at $-0.1$)

\item[(c)]~\\

DC gain:

\[
G(0)=\frac{1}{0.5}=2
\]

Time constant:

\[
\tau = \frac{A}{1/R} = AR = 10
\]

\item[(d)]~\\

First-order low-pass system

Slope: -20 dB/dec

\item[(e)]~\\

PI controller closed-loop:

\[
L(s)=\frac{K_P s+K_I}{s(5s+0.5)}
\]

Characteristic equation:

\[
s(5s+0.5)+K_P s+K_I=0
\]

\[
5s^2+(0.5+K_P)s+K_I=0
\]

Steady-state error:

\[
e_{ss}=0
\]

for stable system

\item[(f)]~\\

PI is used because:

\begin{itemize}
  \item process is integrating-like or slow
  \item steady-state offset unacceptable
  \item P control alone leaves offset
\end{itemize}

\end{enumerate}

\section{Mass-Spring-Damper System}

A mechanical system consists of a mass $m$, damper $c$, and spring $k$.

\[
m\ddot{x}(t) + c\dot{x}(t) + kx(t) = F(t)
\]

Parameters:

\[
m=2,\quad c=4,\quad k=6
\]

Assume zero initial conditions.

\begin{enumerate}

\item[(a)] Derive the transfer function
\[
G(s)=\frac{X(s)}{F(s)}.
\]

\item[(b)] Determine system order and stability.

\item[(c)] Step response analysis:
\begin{enumerate}
\item[(i)] steady-state displacement
\item[(ii)] settling time estimate
\end{enumerate}

\item[(d)] Frequency response interpretation:
\begin{enumerate}
\item[(i)] system type (low/high-pass)
\item[(ii)] DC gain
\item[(iii)] qualitative high-frequency behavior
\end{enumerate}

\item[(e)] Closed-loop with proportional controller $C(s)=K_P$:
\begin{enumerate}
\item[(i)] Derive the closed-loop transfer function.
\item[(ii)] Determine the characteristic equation.
\item[(iii)] Determine the stability condition for $K_P$.
\end{enumerate}

\item[(f)] PI controller:
\[
C(s)=K_P+\frac{K_I}{s}
\]

\begin{enumerate}
\item[(i)] Derive the closed-loop characteristic equation.
\item[(ii)] Determine the steady-state error for a step input.
\end{enumerate}

\end{enumerate}

\ifshowsolutions
Solution:

\begin{enumerate}

\item[(a)]~\\

Laplace transform:

\[
ms^2X(s)+csX(s)+kX(s)=F(s)
\]

\[
G(s)=\frac{X(s)}{F(s)}=\frac{1}{ms^2+cs+k}
\]

Insert values:

\[
G(s)=\frac{1}{2s^2+4s+6}
\]

\item[(b)]~\\

Order: 2

Stability: all coefficients positive → stable

\item[(c)]~\\

Steady-state (step input):

\[
G(0)=\frac{1}{6}
\]

\[
x(\infty)=\frac{1}{6}
\]

Settling time:

dominant pole approx → order 2 stable system:

\[
T_s \approx \frac{4}{\zeta\omega_n}
\]

(accept qualitative estimate)

\item[(d)]~\\

Type: low-pass system

DC gain:

\[
G(0)=\frac{1}{6}
\]

High frequency:

\[
|G(j\omega)| \sim \frac{1}{\omega^2}
\]

→ slope -40 dB/dec

\item[(e)]~\\

Closed loop:

\[
T(s)=\frac{K_P}{2s^2+4s+6+K_P}
\]

Characteristic equation:

\[
2s^2+4s+6+K_P=0
\]

Stability:

all coefficients positive → $K_P>-6$

(practically $K_P>0$)

\item[(f)]~\\

PI controller:

\[
C(s)=\frac{K_P s+K_I}{s}
\]

Characteristic equation:

\[
s(2s^2+4s+6)+K_P s+K_I=0
\]

\[
2s^3+4s^2+(6+K_P)s+K_I=0
\]

Steady-state error:

PI introduces integrator → for step input:

\[
e_{ss}=0
\]

\end{enumerate}

\fi


\section{DC Motor Control System}

A separately excited DC motor is described by the following physical relations.

Electrical subsystem:

\[
u(t) = L\frac{di(t)}{dt} + Ri(t) + e(t)
\]

Back electromotive force:

\[
e(t) = K\,\omega(t)
\]

Mechanical subsystem:

\[
J\frac{d\omega(t)}{dt} + b\,\omega(t) = K\,i(t)
\]

The parameters are:

\[
J=0.02,\quad b=0.2,\quad L=0.5,\quad R=1,\quad K=0.1
\]

Assume zero initial conditions.

\begin{enumerate}

\item[(a)] \textbf{Modeling}~\\

Derive the transfer function

\[
G(s)=\frac{\Omega(s)}{U(s)}
\]

starting from the physical equations.

\item[(b)] \textbf{System representation}~\\

Insert numerical values and express the transfer function in polynomial form.

Determine the system order.

Is the system stable? Justify your answer.

\item[(c)] \textbf{Time-domain behavior}~\\

The system is excited by a unit step input.

\begin{enumerate}
\item[(i)] Determine the steady-state value of the angular velocity.
\item[(ii)] Estimate the settling time using a first-order approximation.
\end{enumerate}

\item[(d)] \textbf{Frequency-domain interpretation}~\\

Consider the transfer function obtained in (b).

\begin{enumerate}
\item[(i)] What type of system (low-pass, high-pass, etc.) is this?
\item[(ii)] What is the DC gain?
\item[(iii)] Explain qualitatively how the magnitude behaves for low and high frequencies.
\end{enumerate}

\item[(e)] \textbf{Proportional control}~\\

The motor is controlled in unity feedback with a proportional controller

\[
C(s)=K_P.
\]

\begin{enumerate}
\item[(i)] Derive the closed-loop transfer function.
\item[(ii)] Determine the characteristic equation.
\item[(iii)] Determine the stability condition for $K_P$.
\end{enumerate}

\item[(f)] \textbf{PI control}~\\

Now consider a PI controller:

\[
C(s)=K_P + \frac{K_I}{s}.
\]

\begin{enumerate}
\item[(i)] Derive the closed-loop characteristic equation.
\item[(ii)] Explain the effect of the integral term on steady-state error.
\item[(iii)] Determine the steady-state error for a step input (qualitatively or using final value theorem).
\end{enumerate}

\end{enumerate}

\ifshowsolutions
Solution:

\begin{enumerate}

\item[(a)] \textbf{Modeling}~\\

Take the Laplace transform under zero initial conditions.

Electrical equation:

\[
U(s)=LsI(s)+RI(s)+E(s)
\]

Using

\[
E(s)=K\Omega(s)
\]

gives

\[
U(s)=(Ls+R)I(s)+K\Omega(s)
\]

Mechanical equation:

\[
Js\Omega(s)+b\Omega(s)=KI(s)
\]

Thus

\[
I(s)=\frac{Js+b}{K}\Omega(s)
\]

Substitute into the electrical equation:

\[
U(s)
=
(Ls+R)\frac{Js+b}{K}\Omega(s)
+
K\Omega(s)
\]

Factor out $\Omega(s)$:

\[
U(s)
=
\left(
\frac{(Ls+R)(Js+b)}{K}
+
K
\right)
\Omega(s)
\]

Therefore,

\[
\frac{\Omega(s)}{U(s)}
=
\frac{K}
{(Js+b)(Ls+R)+K^2}
\]

Hence,

\[
G(s)
=
\frac{K}
{(Js+b)(Ls+R)+K^2}
\]

\item[(b)] \textbf{System representation}~\\

Expand the denominator:

\[
(Js+b)(Ls+R)+K^2
\]

\[
=
JLs^2+(JR+bL)s+bR+K^2
\]

Insert values:

\[
JL=0.02\cdot0.5=0.01
\]

\[
JR=0.02\cdot1=0.02
\]

\[
bL=0.2\cdot0.5=0.10
\]

\[
bR=0.2\cdot1=0.20
\]

\[
K^2=0.01
\]

Therefore,

\[
G(s)
=
\frac{0.1}
{0.01s^2+0.12s+0.21}
\]

The highest power of $s$ is two.

Therefore, the system is a second-order system.

\paragraph{(d) Stability}~\\

The denominator polynomial is

\[
0.01s^2+0.12s+0.21
\]

All coefficients are positive.

For a second-order polynomial

\[
a_2s^2+a_1s+a_0
\]

the system is stable (all poles in the left half-plane) if

\[
a_2>0,\quad a_1>0,\quad a_0>0
\]

Since this condition is fulfilled, the system is stable.

\item[(c)] \textbf{Time-domain behavior}~\\

Step input steady-state value:

\[
G(0)=\frac{0.1}{0.21}\approx 0.476
\]

Thus:

\[
\omega(\infty)\approx 0.476
\]

Settling time approximation:

Normalize denominator:

\[
0.01s^2+0.12s+0.21
\]

Dominant time constant approximation:

\[
\tau \approx \frac{1}{\text{dominant pole}} \approx \frac{1}{2} \approx 0.5 \text{ to } 1 \text{ s}
\]

More robust engineering estimate:

\[
T_s \approx \frac{4}{\text{dominant pole}} \approx 2 \text{ to } 4 \text{ s}
\]

(acceptable qualitative result expected)

\item[(d)] \textbf{Frequency-domain interpretation}~\\

Type of system:

Second-order system with no zeros → low-pass behavior.

DC gain:

\[
G(0)=\frac{0.1}{0.21}\approx 0.476
\]

Qualitative behavior:

\begin{itemize}
  \item For low frequencies: magnitude approximately constant (DC gain)
  \item For high frequencies: denominator dominates → magnitude decreases
  \item Asymptotic slope: approximately \(-40\) dB/decade
\end{itemize}

\item[(e)] \textbf{Proportional control}~\\

Controller:

\[
C(s)=K_P
\]

Closed-loop transfer function:

\[
T(s)=\frac{K_P G(s)}{1+K_P G(s)}
\]

Insert plant:

\[
T(s)=\frac{K_P \cdot 0.1}{0.01s^2+0.12s+0.21+0.1K_P}
\]

Characteristic equation:

\[
0.01s^2+0.12s+0.21+0.1K_P=0
\]

Stability condition:

Second-order system stable if all coefficients positive:

\[
0.01>0,\quad 0.12>0,\quad 0.21+0.1K_P>0
\]

Thus:

\[
K_P>-2.1
\]

For practical control (positive gain):

\[
K_P>0
\]

\item[(f)] \textbf{PI control}~\\

Controller:

\[
C(s)=K_P+\frac{K_I}{s}=\frac{K_P s+K_I}{s}
\]

Open-loop transfer function:

\[
L(s)=\frac{K_P s+K_I}{s}\cdot \frac{0.1}{0.01s^2+0.12s+0.21}
\]

(i) Characteristic equation:

\[
1+L(s)=0
\]

Multiply through:

\[
s(0.01s^2+0.12s+0.21)+0.1(K_P s+K_I)=0
\]

Expand:

\[
0.01s^3+0.12s^2+0.21s+0.1K_P s+0.1K_I=0
\]

Thus:

\[
0.01s^3+0.12s^2+(0.21+0.1K_P)s+0.1K_I=0
\]

(ii) Effect of integral term:

The integral term introduces an additional pole at the origin in the open loop.

Consequences:

\begin{itemize}
  \item eliminates steady-state error for step inputs
  \item increases system type by 1
  \item may reduce stability margin
\end{itemize}

(iii) Steady-state error (step input):

Using final value theorem:

\[
e_{ss}=\lim_{s\to 0} \frac{1}{1+L(s)}
\]

For PI controller:

\[
\lim_{s\to 0} L(s)=\infty
\]

due to \(1/s\) term.

Thus:

\[
e_{ss}=0
\]

for a stable closed-loop system.

\end{enumerate}

\fi

\end{document}